% =============================================================================
% SYSTEM INSTRUCTION TO LLM:
% This file is a PROMPT TEMPLATE.
% You MUST generate a FULL, FINAL, COMPILABLE LaTeX presentation
% based on this template.
%
% CRITICAL RULES:
% 1. DO NOT change anything in the Preamble.
% 2. You MUST delete all [STYLE REF] frames and the AI INSTRUCTION frame.
% 3. You MUST replace all placeholder texts with real content.
% 4. You MUST preserve the collegeBeamer theme and visual style.
% 5. PAGE BUDGET: the final PDF must total 12--18 pages.
%    Scaffolding = title (1) + outline (1) + Q&A (1), and the theme
%    AUTO-INSERTS one divider page at every \section -- never build divider
%    frames by hand. Plan 3--5 sections with 2--4 content frames each.
%    Keep each content slide's body text to roughly 8--10 lines.
% 6. NEVER use \includegraphics. No external image files exist. Draw every chart,
%    diagram and table with tikz / pgfplots / booktabs instead.
% 7. Only these packages are available (already loaded in the preamble): collegebeamer,
%    fontspec, newpxmath, tikz, pgfplots, amsmath, amssymb, pifont, booktabs,
%    colortbl, multirow. Do NOT add \usepackage lines.
% 8. In pgfplots, every \addplot coordinate MUST reuse exactly the labels declared
%    in `symbolic x coords` (any mismatch is a compile error).
% =============================================================================

\documentclass[aspectratio=169,10pt]{beamer}

% -----------------------------------------------------------------------------
% 1. Preamble & Theme (DO NOT CHANGE)
% -----------------------------------------------------------------------------
\usepackage[en,bjtu]{collegebeamer}

\usepackage{tikz}
\usetikzlibrary{positioning}
\usepackage{pgfplots}
\pgfplotsset{compat=1.17}
\usepackage{amsmath,amssymb}
\usepackage{pifont}
\usepackage{booktabs}
\usepackage{colortbl}
\usepackage{multirow}

% --- Font Setup ---
% Text: TeX Gyre Pagella (Palatino clone with real bold/italic/small caps under
% XeLaTeX, loaded from the bundled font/ directory).  Math: newpxmath
% (Palatino-matched math letters).
\usepackage{fontspec}
\setmainfont{texgyrepagella}[
    Path = font/,
    Extension = .otf,
    UprightFont = *-regular,
    BoldFont = *-bold,
    ItalicFont = *-italic,
    BoldItalicFont = *-bolditalic,
]
\let\mathbf\relax % avoid a redefinition clash between bm/beamer and newpxmath
\usepackage{newpxmath}

% NOTE: the theme auto-inserts a highlighted-TOC divider page at every
% \section, and prints "<number> <section name>" under each frame title
% (suppressed automatically before the first \section). Do not redefine these.

% Title page: skips subtitle/author/institute lines cleanly when left empty.
\makeatletter
\def\academic@emptyauthor{\def\inst{\beamer@insttitle}\def\and{\beamer@andtitle}}
\setbeamertemplate{title page}{
    \hspace{-12mm}
    \begin{beamercolorbox}[wd=0.72\textwidth, sep=10pt, leftskip=8mm]{title}
        {\color{maincolor}\rule{2cm}{0.6pt}\par}
        \vskip10pt
        {\usebeamerfont{title}\inserttitle\par}
        \ifx\insertsubtitle\@empty\else
            \vskip8pt{\usebeamerfont{subtitle}\insertsubtitle\par}\fi
        \ifx\insertauthor\academic@emptyauthor\else
            \vskip14pt{\usebeamerfont{author}\usebeamercolor[fg]{author}\insertauthor\par}\fi
        \ifx\insertinstitute\@empty\else
            \vskip4pt{\usebeamerfont{institute}\usebeamercolor[fg]{author}\insertinstitute\par}\fi
        \vskip10pt{\usebeamerfont{date}\usebeamercolor[fg]{date}\insertdate\par}
    \end{beamercolorbox}
}
\makeatother

% -----------------------------------------------------------------------------
% 2. Metadata (LLM MUST FILL)
% -----------------------------------------------------------------------------
\title[Short Title]{Full Presentation Title}
% Caller fills these three fields. Leave empty when unknown; never invent placeholders.
\subtitle{}
\author{}
\institute{}
\date{\today}

\begin{document}

% =============================================================================
% FRAME 1: TITLE PAGE
% =============================================================================
\maketitle

% =============================================================================
% FRAME 2: TABLE OF CONTENTS (auto-generated -- never hardcode section numbers)
% =============================================================================
\begin{frame}{Outline}
    \vspace{0.5em}
    \large
    \tableofcontents
\end{frame}

% =============================================================================
% SECTION 1 (EXAMPLE)
% IMPORTANT: declare \section{...} BEFORE the content frames of each part.
% The theme then automatically inserts a full-color divider page with the
% current section highlighted -- do NOT build any manual transition frames.
% =============================================================================
\section{Title of Part One}

% =============================================================================
% CONTENT FRAME TEMPLATE (LLM MUST USE THIS STRUCTURE)
% =============================================================================
% Content page example 1: Text-left, figure-right layout
\begin{frame}{Content Example: Text Left, Figure Right}
\framesubtitle{} % keep empty; framesubtitle template will show section info automatically
\small
\begin{columns}[T]
\column{0.55\textwidth}
Write the core analysis content on the left side here.
This area is used for background description, problem definition,
logical reasoning, or methodology explanation.
Text should be concise and conclusion-oriented; avoid vague descriptions.

\vspace{0.6em}
\textbf{Key Takeaway:}
State the central message of this slide in one sentence.

\column{0.40\textwidth}
\centering
% Side panel: FEW categories only (4--6 short labels). Compact width. No xlabel.
% For 12-month / many-label series use the full-width chart on the next example page.
% NOTE: keep `enlarge x limits` AFTER the axis-line keys (they reset it),
% and every coordinate label must match `symbolic x coords` exactly.
\begin{tikzpicture}
\begin{axis}[
    width=5.4cm,
    height=3.6cm,
    ybar,
    bar width=11pt,
    ymin=0, ymax=100,
    ytick distance=25,
    symbolic x coords={Q1,Q2,Q3,Q4},
    xtick=data,
    ylabel={Sample index},
    tick label style={font=\scriptsize},
    label style={font=\scriptsize},
    axis y line=left,
    axis x line=bottom,
    axis line style={maincolor!70, -},
    enlarge x limits={abs=0.65cm},
    grid=major,
    grid style={maincolor!12},
]
\addplot[fill=maincolor!40, draw=maincolor] coordinates {
    (Q1,42) (Q2,55) (Q3,61) (Q4,78)
};
\end{axis}
\end{tikzpicture}
\end{columns}
\end{frame}

% Content page example 2: Text-top, figure-bottom layout
\begin{frame}{Content Example: Text Top, Figure Bottom}
\framesubtitle{} % automatically inherits section info
\small

% Upper text area
When a series has many categories (for example 12 months), use this
full-width layout below -- not the narrow side-panel chart.
Keep the text concise so the wide chart remains readable.

\vspace{1em}

% Lower figure: MANY categories (12 months) → MUST use this full-width layout.
% Short labels (<=3 chars) stay horizontal; add `rotate=30, anchor=east` to
% `x tick label style` only when labels are long. Keep `enlarge x limits`
% AFTER the axis-line keys, and match `symbolic x coords` labels exactly.
{\centering
\begin{tikzpicture}
\begin{axis}[
    width=11.5cm, height=4.2cm,
    x tick label style={font=\scriptsize},
    ylabel style={font=\scriptsize},
    xlabel style={font=\scriptsize, yshift=-2pt},
    ymin=1900, ymax=2900,
    symbolic x coords={Jan,Feb,Mar,Apr,May,Jun,Jul,Aug,Sep,Oct,Nov,Dec},
    xtick=data,
    ylabel={Sample level},
    xlabel={Month (illustrative)},
    legend style={at={(0.02,0.98)}, anchor=north west, font=\scriptsize, draw=none, fill=none},
    yticklabel style={font=\scriptsize},
    axis y line=left,
    axis x line=bottom,
    axis line style={maincolor!70, -},
    enlarge x limits=0.04,
    grid=major,
    grid style={maincolor!12},
]
\addplot[
    thick,
    mark=*,
    maincolor,
    mark options={fill=white, draw=maincolor, scale=0.7},
] coordinates {
    (Jan,2040) (Feb,2085) (Mar,2160) (Apr,2330) (May,2345) (Jun,2320)
    (Jul,2420) (Aug,2505) (Sep,2630) (Oct,2740) (Nov,2655) (Dec,2620)
};
\addplot[
    thick,
    dashed,
    maincolor!55,
    mark=square*,
    mark options={fill=white, draw=maincolor!55, scale=0.55},
] coordinates {
    (Jan,2010) (Feb,2050) (Mar,2120) (Apr,2280) (May,2300) (Jun,2290)
    (Jul,2380) (Aug,2460) (Sep,2580) (Oct,2700) (Nov,2620) (Dec,2590)
};
\legend{Series A (sample), Series B (sample)}
\end{axis}
\end{tikzpicture}\par}

\vspace{0.2em}
\hspace*{1.9cm}{\scriptsize\color{darkgray!80}\emph{Source: data compiled from the source document.}}
\end{frame}


% Content page example 3: Figure-left, text-right layout
\begin{frame}{Content Example: Figure Left, Text Right}
\framesubtitle{} 
\small
\begin{columns}[c] % [c] = vertically centered; use [T] for top alignment if text is long

% Left: real TikZ structure diagram (layered pipeline)
\column{0.45\textwidth}
    \centering
    \begin{tikzpicture}[
        font=\scriptsize,
        box/.style={
            draw=maincolor, fill=maincolor!12, rounded corners=2pt,
            minimum height=0.62cm, text width=1.35cm, align=center
        },
        layer title/.style={font=\scriptsize\bfseries, text=maincolor}
    ]
        % Data layer (panel top leaves reserved space for the caption)
        \fill[maincolor!6, rounded corners=3pt] (-0.25,1.05) rectangle (4.85,2.60);
        \node[layer title, anchor=north west] at (-0.15,2.55) {Data layer};
        \node[box] (s1) at (0.7,1.50) {Sensors};
        \node[box] (s2) at (2.5,1.50) {Ingest};
        \node[box] (s3) at (4.2,1.50) {Store};
        % Decision layer
        \fill[maincolor!6, rounded corners=3pt] (-0.25,-1.25) rectangle (4.85,0.70);
        \node[layer title, anchor=north west] at (-0.15,0.65) {Decision layer};
        \node[box] (p1) at (0.7,-0.30) {Feature};
        \node[box] (p2) at (2.5,-0.30) {Model};
        \node[box] (p3) at (4.2,-0.30) {Act};
        % Edges (route cross-layer links orthogonally around the panels so
        % they never slice through captions or other boxes)
        \draw[->, thick, maincolor] (s1) -- (s2);
        \draw[->, thick, maincolor] (s2) -- (s3);
        \draw[->, thick, maincolor, rounded corners=2pt]
            (s2.south) -- (2.5,0.875) -- (-0.65,0.875) -- (-0.65,-0.30) -- (p1.west);
        \draw[->, thick, maincolor] (p1) -- (p2);
        \draw[->, thick, maincolor] (p2) -- (p3);
        \draw[->, thick, maincolor, dashed]
            (p3.south) -- ++(0,-0.30) -| (p1.south)
            node[pos=0.25, below, font=\tiny, text=maincolor] {feedback};
    \end{tikzpicture}

% Right: text analysis area (50\% width)
\column{0.50\textwidth}
    \textbf{Figure / Object Analysis:}
    
    This section provides a detailed interpretation of the visual element
    on the left. Suitable for presenting:
    \begin{itemize}
        \item Physical structure of experimental apparatus
        \item Specific module details of an algorithm or model
        \item Domain-specific data analysis
    \end{itemize}

    \vspace{1em}
    \textbf{Key Findings:}
    \begin{enumerate}
        \item \ctextbf{Observation 1}: The figure on the left shows\ldots
        \item \ctextbf{Observation 2}: Compared to expectations, the results exhibit\ldots
    \end{enumerate}
\end{columns}
\end{frame}

% Content page example 4: Definition block + displayed equation
% Use this pattern whenever the source document is math-heavy.
\begin{frame}{Content Example: Definition and Key Equation}
\framesubtitle{}
\small
\begin{block}{Definition (state the core assumption or theorem here)}
    Summarize the central definition, assumption, or theorem of this part
    in one or two sentences.
\end{block}

\vspace{0.6em}
The key relation then follows directly:
\begin{equation*}
    \hat{y}_t = f_\theta(x_t) + \varepsilon_t,
    \qquad
    \mathcal{L}(\theta) = \frac{1}{T}\sum_{t=1}^{T}\bigl(y_t - \hat{y}_t\bigr)^2 .
\end{equation*}

\vspace{0.4em}
\begin{itemize}
    \item Explain every symbol in one line directly below the equation.
    \item Display only the one or two equations that carry the argument;
          keep everything else inline.
\end{itemize}
\end{frame}

% =============================================================================
% [STYLE REF] TABLE PAGE (DELETE IN FINAL OUTPUT)
% =============================================================================
\begin{frame}{[STYLE REF] Table Style Example}
\framesubtitle{}
\small
\centering
\renewcommand{\arraystretch}{1.3}
\footnotesize
% Header: maincolor fill + white bold text; body: plain rows + \bottomrule.
% Use \rowcolors zebra striping ONLY for tables with 6+ rows, and then add
% \noalign{\vskip2pt} before \bottomrule to keep the stripe off the rule.
\begin{tabular}{p{0.2\textwidth} c c c}
\rowcolor{maincolor}
\textcolor{white}{\textbf{Metric}} & \textcolor{white}{\textbf{A}} &
\textcolor{white}{\textbf{B}} & \textcolor{white}{\textbf{Conclusion}} \\
Item 1 & 12 & 15 & Room for improvement \\
Item 2 & 20 & 18 & Stable \\
\bottomrule
\end{tabular}
\end{frame}

% =============================================================================
% [STYLE REF] TIKZ FLOW PAGE (DELETE IN FINAL OUTPUT)
% =============================================================================
\begin{frame}{[STYLE REF] Flowchart Style Example}
\framesubtitle{}
\centering
\begin{tikzpicture}[
node distance=0.9cm and 0.7cm,
process/.style={rectangle,draw=maincolor,fill=maincolor!10,
minimum height=0.95cm,text width=1.9cm,align=center}
]
\node[process] (a) {Collect};
\node[process,right=of a] (b) {Analyze};
\node[process,right=of b] (c) {Decide};
\node[process,below=of b] (d) {Revise};
\draw[->,thick,maincolor] (a)--(b);
\draw[->,thick,maincolor] (b)--(c);
\draw[->,thick,maincolor] (b)--(d)
    node[midway,right,font=\tiny,text=maincolor] {gap};
\draw[->,thick,maincolor,dashed] (d.west) -| (a.south)
    node[near end,left,font=\tiny,text=maincolor] {retry};
\end{tikzpicture}
\end{frame}

% =============================================================================
% [STYLE REF] TWO-COLUMN ANALYSIS (DELETE IN FINAL OUTPUT)
% =============================================================================
\begin{frame}{[STYLE REF] Two-Column Analysis Example}
\framesubtitle{}
\small
\begin{columns}
\column{0.52\textwidth}
Left column for argumentation, definitions, and derivations.

\column{0.43\textwidth}
Right column for action items, illustrations, or comparisons.
\end{columns}
\end{frame}

% =============================================================================
% AI INSTRUCTION PAGE (DELETE IN FINAL OUTPUT)
% =============================================================================
\begin{frame}[allowframebreaks]{AI GENERATION INSTRUCTIONS}
\framesubtitle{}
You MUST:
\begin{itemize}
\item Delete all [STYLE REF] frames and this instruction frame
\item Replace all placeholders with real content
\item Declare \texttt{\textbackslash section\{...\}} before the content frames of each
      part; the theme auto-inserts a highlighted divider page at every
      \texttt{\textbackslash section} -- NEVER build manual transition or outline frames
\item PAGE BUDGET: 12--18 PDF pages total = title + outline + Q\&A (3 pages)
      + one auto divider per section + 2--4 content frames per section,
      for 3--5 sections
\item Lay out every content page with one of the demonstrated example layouts;
      full-width (no columns) is REQUIRED for charts with more than 6 categories
\item In every chart, each \texttt{\textbackslash addplot} coordinate must reuse exactly
      the labels declared in \texttt{symbolic x coords}, and keep
      \texttt{enlarge x limits} AFTER the axis-line keys
\item Keep information density high (no empty slides)
\item Output FULL LaTeX code only
\end{itemize}
\end{frame}

% =============================================================================
% FINAL PAGE
% =============================================================================
\QApage

\end{document}
